\section*{Declarations}
\paragraph{Availability of data and materials} The versioned Figshare v3 record \citep{Abaee2026code} contains the 25-check exact-integer grid verifier, results/report, three figure scripts and their figures, and preserved search records; its attached manuscript is Paper1 v42, not this Paper06 v75. The SafeTransition 1.3.0 Figshare v1 deposit \citep{Abaee2026software} contains the software library, tests and benchmark artifacts, and a Paper1 verification artifact; it does not contain this Paper06 manuscript. The Lean formalisation is pinned at commit \texttt{6f626ee48b2a3a59282b5df5adda76c8d5993fdc}: \href{https://github.com/MIKEAA2020/general-sustainability/blob/6f626ee48b2a3a59282b5df5adda76c8d5993fdc/lean/Formalizations/P1_AssessmentSeparation.lean}{source}. Scripts/logs for the ladder, off-diagonal, continuum and common-shock check lists described in SI~S7 were not present in the two inspected Figshare packages or the inspected v71 source bundle; those checks are reported in SI~S7, not claimed to be reproducible from these DOI packages alone. The present v75 article and SI source/PDF are versioned separately in the companion archive. Published DFO biomass/LRP values are used only as a limited scale/context anchor: no parameter is estimated, no new empirical dataset is generated or fitted, and no empirical fit, forecast, or validation is claimed.
\paragraph{Competing interests} The author declares no competing interests.
\paragraph{Funding} Not applicable.
\paragraph{Author contributions} A.A.\ conceptualized the entire work, wrote, reviewed and edited the manuscript.
\paragraph{AI declaration} GLM (Z.ai), Qwen (Alibaba Cloud) and DeepSeek AI assisted with drafting and iterative review; all results, proofs and verifications were checked by the author.

\section*{Supplementary Information}
The Supplementary Information (SI) contains: S1 abbreviations; S2 remaining proofs (persistence and erasure, the substitutability ladder, the common-shock verdict); S3 the substitutability spectrum in full (critical-elasticity landscape, off-diagonal criterion, fixed-sum structure, continuum closure, the Living Planet Index identification); S4 the common-shock variant in full (coupled and replace readings, depth family); S5 the resource-transition benchmark in full (schedule, dashboard arithmetic, fisheries translation table, limited scale/context anchor); S6 practitioner material (terminology map, audit checklist, minimum report, data requirements); S7 the verification record; S8 positioning notes per literature; S9 the framework extensions in full (extended definitions, governance constructors and the commons obstruction, the intergenerational and composition extensions, the planetary-boundaries note, declared conjectures), at the status each item carries. File: \texttt{paper06\_assessment\_separation\_v75\_supplementary} (\texttt{.tex}/\texttt{.pdf}).

\begin{thebibliography}{99}
\makeatletter\renewcommand{\@biblabel}[1]{}\makeatother

\bibitem[Abaee(2026a)]{Abaee2026obs} Abaee, A. (2026a). An obstruction calculus for viability under incomplete observation. Manuscript; Zenodo. \url{https://doi.org/10.5281/zenodo.22552616}.

\bibitem[Abaee(2026b)]{Abaee2026ecomod} Abaee, A. (2026b). \emph{A rising biocapacity total can mask ecological-capital land drawdown: a two-land model}. Companion manuscript, c10. Exact source and PDF pinned to commit \texttt{6f626ee48b2a3a59282b5df5adda76c8d5993fdc}: \href{https://github.com/MIKEAA2020/general-sustainability/blob/6f626ee48b2a3a59282b5df5adda76c8d5993fdc/agent%202%20productivity%20illusion/output_v45_ecomod_continued_10_gemini38_v4base_si_review_20261026/ecomod_continued_10.pdf}{PDF}; \href{https://github.com/MIKEAA2020/general-sustainability/blob/6f626ee48b2a3a59282b5df5adda76c8d5993fdc/agent%202%20productivity%20illusion/output_v45_ecomod_continued_10_gemini38_v4base_si_review_20261026/ecomod_continued_10.tex}{source}.

\bibitem[Abaee(2026c)]{Abaee2026cod} Abaee, A. (2026c). Does a surplus-production ladder improve forecasts of Northern cod? A scored test on NAFO 2J3KL. Zenodo. \url{https://doi.org/10.5281/zenodo.22553609}.

\bibitem[Abaee(2026d)]{Abaee2026code} Abaee, A. (2026d). Verification code and figure pipeline for ``Aggregate indices and transition safety''. Figshare, version 3. The attached manuscript is Paper1 v42, an earlier edition in the same research line, not this Paper06 v75 manuscript. \url{https://doi.org/10.6084/m9.figshare.33764023.v3}.
\bibitem[Abaee(2026e)]{Abaee2026software} Abaee, A. (2026e). \emph{SafeTransition: exact rational certification of transition safety for sustainability assessment}. Software release 1.3.0, Figshare, version 1. \url{https://doi.org/10.6084/m9.figshare.33971389.v1}.

\bibitem[Asheim(1994)]{Asheim1994} Asheim, G.~B. (1994). Net national product as an indicator of sustainability. \emph{Scandinavian Journal of Economics}, 96(2), 257--265.

\bibitem[Athan and Papalambros(1996)]{AthanPapalambros1996} Athan, T. W., Papalambros, P. Y. (1996). A note on weighted criteria methods for compromise solutions in multi-objective optimization. \emph{Engineering Optimization}, 27(2), 155--176.

\bibitem[Aubin and Catt\'e(2002)]{AubinCatte2002} Aubin, J.-P., and Catt\'e, F. (2002). Bilateral fixed-points and algebraic properties of viability kernels and capture basins of sets. \emph{Set-Valued Analysis}, 10(4), 379--416.

\bibitem[Aubin et~al.(2011)]{Aubin2011} Aubin, J.-P., Bayen, A.~M., and Saint-Pierre, P. (2011). \emph{Viability Theory: New Directions}, 2nd ed. Birkh\"auser, Boston.

\bibitem[Aubin(1991)]{Aubin1991} Aubin, J.-P. (1991). \emph{Viability Theory}. Birkh\"auser, Boston.

\bibitem[Ayres et al.(2001)]{Ayres2001} Ayres, R. U., van den Bergh, J. C. J. M., Gowdy, J. M. (2001). Strong versus weak sustainability: economics, natural sciences, and ``consilience''. \emph{Environmental Ethics}, 23(2), 155--168.

\bibitem[Becker et~al.(2017)]{Becker2017} Becker, W., Saisana, M., Paruolo, P., and Vandecasteele, I. (2017). Weights and importance in composite indicators: Closing the gap. \emph{Ecological Indicators}, 80, 12--22.

\bibitem[Beckerman(1994)]{Beckerman1994} Beckerman, W. (1994). `Sustainable development': is it a useful concept? \emph{Environmental Values}, 3(3), 191--209.

\bibitem[Ben-Tal et~al.(2004)]{BenTal2004} Ben-Tal, A., Goryashko, A., Guslitzer, E., and Nemirovski, A. (2004). Adjustable robust solutions of uncertain linear programs. \emph{Mathematical Programming}, 99(2), 351--376.

\bibitem[Ben-Tal et al.(2009)]{BenTal2009} Ben-Tal, A., El Ghaoui, L., Nemirovski, A. (2009). \emph{Robust Optimization}. Princeton University Press.

\bibitem[Bertsimas and Goyal(2010)]{BertsimasGoyal2010} Bertsimas, D., Goyal, V. (2010). On the power of robust solutions in two-stage stochastic and adaptive optimization problems. \emph{Mathematics of Operations Research}, 35(2), 284--305.

\bibitem[Boos(2015)]{Boos2015} Boos, A. (2015). Genuine savings as an indicator for ``weak'' sustainability: Critical survey and possible ways forward in measuring weak sustainability. \emph{Sustainability}, 7(4), 4146--4163.

\bibitem[Cairns and Martinet(2014)]{CairnsMartinet2014} Cairns, R.~D., and Martinet, V. (2014). An environmental-economic measure of sustainable development. \emph{European Economic Review}, 69, 4--17.

\bibitem[Cardaliaguet(1996)]{Cardaliaguet1996} Cardaliaguet, P. (1996). A differential game with two players and one target. \emph{SIAM Journal on Control and Optimization}, 34(4), 1441--1460.

\bibitem[Cinelli et~al.(2014)]{Cinelli2014} Cinelli, M., Coles, S.~R., and Kirwan, K. (2014). Analysis of the potentials of multi criteria decision analysis methods to conduct sustainability assessment. \emph{Ecological Indicators}, 46, 138--148.

\bibitem[Daly(1990)]{Daly1990} Daly, H.~E. (1990). Toward some operational principles of sustainable development. \emph{Ecological Economics}, 2(1), 1--6.

\bibitem[Daly(1995)]{Daly1995} Daly, H. E. (1995). On Wilfred Beckerman's critique of sustainable development. \emph{Environmental Values}, 4(1), 49--55.

\bibitem[Das and Dennis(1997)]{DasDennis1997} Das, I., and Dennis, J.~E. (1997). A closer look at drawbacks of minimizing weighted sums of objectives for Pareto set generation in multicriteria optimization problems. \emph{Structural Optimization}, 14, 63--69.

\bibitem[Dasgupta and M\"aler(2000)]{DasguptaMaler2000} Dasgupta, P., and M\"aler, K.-G. (2000). Net national product, wealth, and social well-being. \emph{Environment and Development Economics}, 5(1), 69--93.

\bibitem[De Lara and Doyen(2008)]{DeLaraDoyen2008} De Lara, M., and Doyen, L. (2008). \emph{Sustainable Management of Natural Resources: Mathematical Models and Methods}. Springer, Berlin.

\bibitem[Delage et al.(2019)]{Delage2019} Delage, E., Kuhn, D., Wiesemann, W. (2019). ``Dice''-sion-making under uncertainty: when can a random decision reduce risk? \emph{Management Science}, 65(7), 3282--3301.

\bibitem[DFO(2016)]{DFO2016} DFO (2016). \emph{Stock Assessment of Northern Cod (NAFO Divs.\ 2J3KL) in 2016}. Canadian Science Advisory Secretariat, Science Advisory Report 2016/026. Fisheries and Oceans Canada, Ottawa.

\bibitem[Dietz and Neumayer(2007)]{DietzNeumayer2007} Dietz, S., Neumayer, E. (2007). Weak and strong sustainability in the SEEA: concepts and measurement. \emph{Ecological Economics}, 61(4), 617--626.

\bibitem[Doyen and Gajardo(2020)]{DoyenGajardo2020} Doyen, L., and Gajardo, P. (2020). Sustainability standards, multicriteria maximin, and viability. \emph{Natural Resource Modeling}, 33(3), e12250.

\bibitem[Doyen et~al.(2012)]{Doyen2012} Doyen, L., Th\'ebaud, O., B\'en\'e, C., Martinet, V., Gourguet, S., Bertignac, M., Fifas, S., and Blanchard, F. (2012). A stochastic viability approach to ecosystem-based fisheries management. \emph{Ecological Economics}, 75, 32--42.

\bibitem[Ehrgott(2005)]{Ehrgott2005} Ehrgott, M. (2005). \emph{Multicriteria Optimization} (2nd ed.). Springer.

\bibitem[Ekins et~al.(2003)]{Ekins2003} Ekins, P., Simon, S., Deutsch, L., Folke, C., and De Groot, R. (2003). A framework for the practical application of the concepts of critical natural capital and strong sustainability. \emph{Ecological Economics}, 44(2), 165--185.

\bibitem[Fabian et~al.(2010)]{Fabian2010} Fabian, M.~J., Henrion, R., Kruger, A.~Y., and Outrata, J.~V. (2010). Error bounds: necessary and sufficient conditions. \emph{Set-Valued and Variational Analysis}, 18(2), 121--149.

\bibitem[Fanning et~al.(2022)]{Fanning2022} Fanning, A.~L., O'Neill, D.~W., Hickel, J., and Roux, N. (2022). The social shortfall and ecological overshoot of nations. \emph{Nature Sustainability}, 5(1), 26--36.

\bibitem[Filippov(1988)]{Filippov1988} Filippov, A.~F. (1988). \emph{Differential Equations with Discontinuous Righthand Sides}. Kluwer, Dordrecht.

\bibitem[Fischer et~al.(2022)]{Fischer2022} Fischer, S.~M., Joy, M.~K., Abrahamse, W., Milfont, T.~L., and Petherick, L.~M. (2022). The use and misuse of composite environmental indices. \emph{bioRxiv}, 2022.03.15.484501. \url{https://doi.org/10.1101/2022.03.15.484501}.

\bibitem[Frankowska(1989)]{Frankowska1989} Frankowska, H. (1989). Optimal trajectories associated with a solution of contingent Hamilton--Jacobi equations. \emph{Applied Mathematics and Optimization}, 19, 291--311.

\bibitem[Gao et~al.(2023)]{Gao2023} Gao, P., Wang, Y., Wang, H., Song, C., Ye, S., and Wang, X. (2023). A Pareto front-based approach for constructing composite index of sustainability without weights: A comparative study of implementations. \emph{Ecological Indicators}, 155, 110919.

\bibitem[Hanley et~al.(1999)]{Hanley1999} Hanley, N., Moffatt, I., Faichney, R., and Wilson, M. (1999). Measuring sustainability: A time series of alternative indicators for Scotland. \emph{Ecological Economics}, 28(1), 55--73.

\bibitem[Hickel(2020)]{Hickel2020} Hickel, J. (2020). The sustainable development index: Measuring the ecological efficiency of human development in the Anthropocene. \emph{Ecological Economics}, 167, 106331.

\bibitem[Keeney and Raiffa(1976)]{KeeneyRaiffa1976} Keeney, R.~L., and Raiffa, H. (1976). \emph{Decisions with Multiple Objectives: Preferences and Value Tradeoffs}. Wiley, New York.

\bibitem[Koski(1985)]{Koski1985} Koski, J. (1985). Defectiveness of weighting method in multicriterion optimization of structures. \emph{Communications in Applied Numerical Methods}, 1(6), 333--337.

\bibitem[Krause et al.(2011)]{KrauseRoperGolovin2011} Krause, A., Roper, A., Golovin, D. (2011). Randomized sensing in adversarial environments. In \emph{Proceedings of the 22nd International Joint Conference on Artificial Intelligence (IJCAI)}, 2133--2139.

\bibitem[Kuhn(1953)]{Kuhn1953} Kuhn, H. W. (1953). Extensive games and the problem of information. In H. W. Kuhn, A. W. Tucker (Eds.), \emph{Contributions to the Theory of Games II} (Annals of Mathematics Studies 28, pp. 193--216). Princeton University Press.

\bibitem[Lade et~al.(2020)]{Lade2020} Lade, S.~J., Steffen, W., de Vries, W., Carpenter, S.~R., Donges, J.~F., Gerten, D., Hoff, H., Newbold, T., Richardson, K., and Rockstr\"om, J. (2020). Human impacts on planetary boundaries amplified by Earth system interactions. \emph{Nature Sustainability}, 3(2), 119--128.

\bibitem[Lygeros et~al.(1999)]{Lygeros1999} Lygeros, J., Tomlin, C., and Sastry, S. (1999). Controllers for reachability specifications for hybrid systems. \emph{Automatica}, 35(3), 349--370.

\bibitem[Main and Randour(2024)]{MainRandour2024} Main, J. C. A., Randour, M. (2024). Different strokes in randomised strategies: revisiting Kuhn's theorem under finite-memory assumptions. \emph{Information and Computation}, 301, 105229.

\bibitem[Martinet(2011)]{Martinet2011} Martinet, V. (2011). A characterization of sustainability with indicators. \emph{Journal of Environmental Economics and Management}, 61(2), 183--197.

\bibitem[Martinez-Alier et~al.(1998)]{MartinezAlier1998} Martinez-Alier, J., Munda, G., and O'Neill, J. (1998). Weak comparability of values as a foundation for ecological economics. \emph{Ecological Economics}, 26(3), 277--286.

\bibitem[Messac and Mattson(2002)]{MessacMattson2002} Messac, A., Mattson, C. A. (2002). Generating well-distributed sets of Pareto points for engineering design using physical programming. \emph{Optimization and Engineering}, 3(4), 431--450.

\bibitem[Miettinen(1999)]{Miettinen1999} Miettinen, K. (1999). \emph{Nonlinear Multiobjective Optimization}. Kluwer.

\bibitem[Munda and Nardo(2009)]{MundaNardo2009} Munda, G., Nardo, M. (2009). Noncompensatory/nonlinear composite indicators for ranking countries: a defensible setting. \emph{Applied Economics}, 41(12), 1513--1523.

\bibitem[Munda(2005)]{Munda2005} Munda, G. (2005). Multiple criteria decision analysis and sustainable development. In Figueira, J., Greco, S., and Ehrgott, M. (eds.), \emph{Multiple Criteria Decision Analysis: State of the Art Surveys}. Springer, New York, 953--986.

\bibitem[Nardo et~al.(2008)]{Nardo2008} Nardo, M., Saisana, M., Saltelli, A., Tarantola, S., Hoffmann, A., and Giovannini, E. (2008). \emph{Handbook on Constructing Composite Indicators: Methodology and User Guide}. OECD Publishing, Paris.

\bibitem[von Neumann(1928)]{vonNeumann1928} von Neumann, J. (1928). Zur Theorie der Gesellschaftsspiele. \emph{Mathematische Annalen}, 100, 295--320.

\bibitem[Neumayer(2013)]{Neumayer2013} Neumayer, E. (2013). \emph{Weak versus Strong Sustainability: Exploring the Limits of Two Opposing Paradigms}, 4th ed. Edward Elgar, Cheltenham.

\bibitem[O'Neill et~al.(2018)]{ONeill2018} O'Neill, D.~W., Fanning, A.~L., Lamb, W.~F., and Steinberger, J.~K. (2018). A good life for all within planetary boundaries. \emph{Nature Sustainability}, 1(2), 88--95.

\bibitem[Pearce and Atkinson(1993)]{PearceAtkinson1993} Pearce, D. W., Atkinson, G. D. (1993). Capital theory and the measurement of sustainable development: an indicator of ``weak'' sustainability. \emph{Ecological Economics}, 8(2), 103--108.

\bibitem[Raworth(2012)]{Raworth2012} Raworth, K. (2012). A safe and just space for humanity: Can we live within the doughnut? \emph{Oxfam Policy and Practice: Climate Change and Resilience}, 8(1), 1--26.

\bibitem[Rockstr\"om et~al.(2009)]{Rockstrom2009} Rockstr\"om, J., Steffen, W., Noone, K., et al. (2009). A safe operating space for humanity. \emph{Nature}, 461, 472--475.

\bibitem[Roy(1996)]{Roy1996} Roy, B. (1996). \emph{Multicriteria Methodology for Decision Aiding}. Kluwer, Dordrecht.

\bibitem[Saint-Pierre(1994)]{SaintPierre1994} Saint-Pierre, P. (1994). Approximation of the viability kernel. \emph{Applied Mathematics and Optimization}, 29, 187--209.

\bibitem[Sch\"ar et~al.(2025)]{Schar2025} Sch\"ar, S., Pohl, E., and Geldermann, J. (2025). Analysing the compensatory properties of the outranking approach PROMETHEE. \emph{Journal of Multi-Criteria Decision Analysis}, 32, e70013.

\bibitem[Schaefer(1954)]{Schaefer1954} Schaefer, M.~B. (1954). Some aspects of the dynamics of populations important to the management of commercial marine fisheries. \emph{Inter-American Tropical Tuna Commission Bulletin}, 1(2), 23--56.

\bibitem[Sessa et al.(2020)]{Sessa2020} Sessa, P. G., Bogunovic, I., Kamgarpour, M., Krause, A. (2020). Mixed strategies for robust optimization of unknown objectives. In \emph{Proceedings of the 23rd International Conference on Artificial Intelligence and Statistics (AISTATS)}, PMLR 108, 2970--2980.

\bibitem[Sinha et al.(2018)]{Sinha2018} Sinha, A., Fang, F., An, B., Kiekintveld, C., Tambe, M. (2018). Stackelberg security games: looking beyond a decade of success. In \emph{Proceedings of the 27th International Joint Conference on Artificial Intelligence (IJCAI)}, 5494--5501.

\bibitem[Sion(1958)]{Sion1958} Sion, M. (1958). On general minimax theorems. \emph{Pacific Journal of Mathematics}, 8(1), 171--176.

\bibitem[Solow(1974)]{Solow1974} Solow, R.~M. (1974). Intergenerational equity and exhaustible resources. \emph{Review of Economic Studies}, 41, 29--45.

\bibitem[Usubiaga-Lia\~no(2025)]{Usubiaga2025} Usubiaga-Lia\~no, A. (2025). Strong sustainability in the SEEA and the wider indicator debate. \emph{One Ecosystem}, 10, e141086.

\bibitem[Vincke(1992)]{Vincke1992} Vincke, P. (1992). \emph{Multicriteria Decision-Aid}. Wiley, Chichester.

\bibitem[Vorobeychik and Li(2014)]{VorobeychikLi2014} Vorobeychik, Y., Li, B. (2014). Optimal randomized classification in adversarial settings. In \emph{Proceedings of the 13th International Conference on Autonomous Agents and Multiagent Systems (AAMAS)}, 485--492.

\bibitem[Warga(1972)]{Warga1972} Warga, J. (1972). \emph{Optimal Control of Differential and Functional Equations}. Academic Press, New York.

\bibitem[Wei and Zhang(2024)]{WeiZhang2024} Wei, N., Zhang, P. (2024). Adjustability in robust linear optimization. \emph{Mathematical Programming}, 208, 581--628.

\bibitem[World Bank(2011)]{WorldBank2011} World Bank (2011). \emph{The Changing Wealth of Nations: Measuring Sustainable Development in the New Millennium}. World Bank, Washington, DC.

\end{thebibliography}

\end{document}
